\section{Geometric phase, Berry/Wilczek--Zee connections, and gauge structure}
\label{sec:geophase}

\subsection{Per-wavepacket band isolation and polarization transport}

The complex-envelope description is applied per band-isolated excitation.
Each wavepacket is treated as narrowband around its own local carrier angular frequency $\omega_0$,
determined by the linearized carrier operator about the slowly varying background at the wavepacket's
location; no globally coherent universe-wide carrier is assumed.
A band-isolated wavepacket can be written locally as
\begin{equation}
  \vecu(x,t)\;\approx\;\mathrm{Re}\Big\{ a(x,t)\,\bm{\epsilon}(x,t)\,e^{\ii\theta(x,t)} \Big\},
  \qquad
  \theta(x,t)=\phi(x,t)-\omega_0 t,
  \label{eq:quasi-monochromatic}
\end{equation}
where $a(x,t)\ge 0$ varies slowly on the envelope scale, $\bm{\epsilon}(x,t)\in\C^{N}$ is a normalized
polarization state ($\bm{\epsilon}^\dagger\bm{\epsilon}=1$), and $\omega_0$ is the wavepacket's local carrier
angular frequency.
The internal dimension $N$ counts the relevant components (e.g.\ embedding-direction components and/or a
coarse-grained mode basis).

A central point is that polarization transport is not unique: the same physical state can be represented in many
local phase conventions.
Geometric phase formalizes this redundancy as a connection on a bundle of polarization states.

\subsection{Rank-1 Berry connection (U(1))}

For a single isolated polarization mode, represent the local state by a normalized vector
$\lvert u(x)\rangle\in\C^{N}$ that depends smoothly on spacetime position $x^\mu$
(here $x^\mu=(c\,t,x^1,x^2,x^3)$ built from the intrinsic coordinates and a chosen characteristic speed $c$ of the
carrier branch).
Define the Berry connection
\begin{equation}
  a_\mu(x) := \ii\,\langle u(x)\,|\,\partial_\mu u(x)\rangle.
  \label{eq:berry-connection}
\end{equation}
Under a local phase choice $\lvert u(x)\rangle\mapsto e^{-\ii\chi(x)}\lvert u(x)\rangle$ one finds
\begin{equation}
  a_\mu \mapsto a_\mu + \partial_\mu \chi,
  \label{eq:berry-gauge}
\end{equation}
so only gauge-invariant quantities carry physical meaning.
The Berry curvature is the field strength
\begin{equation}
  f_{\mu\nu} := \partial_\mu a_\nu - \partial_\nu a_\mu,
  \label{eq:berry-curvature}
\end{equation}
and the holonomy (geometric phase) around a closed loop $\Gamma$ is
\begin{equation}
  \gamma_\Gamma
  = \oint_\Gamma a_\mu\,\dd x^\mu
  = \int_{\Sigma:\,\partial\Sigma=\Gamma} f_{\mu\nu}\,\dd\Sigma^{\mu\nu}.
  \label{eq:berry-holonomy}
\end{equation}

\subsection{What the geometric phase is ``of'': eigenmode bundles and adiabaticity}
\label{subsec:berry-of}

The connection above is not attached to an arbitrary decomposition of a real field into amplitude and phase.
In the present setting one should choose $\lvert u(x)\rangle$ as a (locally defined) eigenvector of the
linearized carrier operator about the slowly varying background configuration at $x$.
As $x$ varies, these eigenvectors form an eigenbundle, and the Berry/Wilczek--Zee connection is the natural
connection on that bundle; this is the same eigenbundle picture that underlies Berry-phase effects in
crystalline electronic systems \citep{XiaoChangNiu2010Berry}, here applied to the polarization eigenbundle of
the mechanical carrier rather than to electronic Bloch bands.

A clean rank-1 ($U(1)$) reduction requires that the dynamics remains in a single isolated band.
If $\Delta\omega_{\mathrm{gap}}(x)$ is the local separation to the nearest competing band and
$\Omega_{\mathrm{mix}}(x)$ is the rate of non-adiabatic mode conversion induced by background gradients,
then the adiabatic condition
\begin{equation}
  \Omega_{\mathrm{mix}}(x)\ll \Delta\omega_{\mathrm{gap}}(x)
  \label{eq:adiabatic-condition}
\end{equation}
should hold on the envelope scale.
When this fails, a measured ``Berry signal'' typically reflects mode mixing or beating, and the correct reduction
is to retain the relevant multi-dimensional subspace (the Wilczek--Zee description).
In the brane model, the pre-stress parameter $\alpha$ and the wavelength regime control branch separation and
therefore how well an Abelian gauge description can apply.

\subsection{Diagnostic protocol and artifact controls (numerical extraction)}
\label{subsec:diagnostic-protocol}

The Berry/Wilczek--Zee connection is meaningful only when it is tied to a specific eigenbundle (or transported
subspace) and when the reported observables are gauge-invariant.
For numerical experiments on the lattice substrate, the following protocol provides a concrete, reproducible
procedure to extract holonomy/curvature and to rule out common artifacts.

\begin{center}
\fbox{\begin{minipage}{0.95\linewidth}
\textbf{Protocol}
\begin{enumerate}
\item \textbf{Carrier isolation / demodulation.}
Choose a carrier frequency $\omega_0$ and isolate a narrowband sector (e.g.\ by band-pass filtering in time or by
short-time Fourier methods).
Work with a complex envelope representation consistent with \Cref{eq:quasi-monochromatic}.

\item \textbf{Define the instantaneous polarization object.}
At each spacetime point (or each sample along a loop), define either
(i) a normalized rank-1 state $\lvert u(x)\rangle$ for an isolated band, or
(ii) an $n$-dimensional subspace projector $P(x)$ (or an orthonormal frame $\mathbf{E}(x)$) for a near-degenerate
sector.
The definition should be tied to the \emph{local linearized carrier problem} about the slowly varying background at $x$
(cf.\ \Cref{subsec:berry-of}).

\item \textbf{Measure the adiabaticity/gap indicators.}
Estimate a local band gap $\Delta\omega_{\mathrm{gap}}(x)$ (separation to the nearest competing band) and a mode-leakage
proxy along the transport path (e.g.\ growth of energy outside the chosen band/subspace).
Use these to empirically validate the adiabatic criterion \Cref{eq:adiabatic-condition}.

\item \textbf{Holonomy from overlaps (gauge-invariant).}
For a discretized loop $\Gamma: x_0\!\to x_1\!\to\cdots\to x_m=x_0$:
\begin{itemize}
\item \emph{Abelian (rank-1):} define links
$U_k=\frac{\langle u(x_k)\,|\,u(x_{k+1})\rangle}{|\langle u(x_k)\,|\,u(x_{k+1})\rangle|}\in U(1)$ and report the loop phase
$\gamma_\Gamma=\arg\!\big(\prod_{k=0}^{m-1}U_k\big)$.
\item \emph{Non-Abelian (rank-$n$):} define overlap matrices $M_k=\mathbf{E}(x_k)^\dagger\mathbf{E}(x_{k+1})$ and extract their
unitary part (e.g.\ by polar decomposition) as $U_k\in U(n)$.
Report the holonomy matrix $U_\Gamma=\prod_{k=0}^{m-1}U_k$ and gauge-invariant summaries such as
$\mathrm{spec}(U_\Gamma)$ and $\mathrm{tr}(U_\Gamma)$.
\end{itemize}

\item \textbf{Curvature by small loops / plaquettes.}
Approximate curvature by evaluating holonomy around small loops and dividing by enclosed area (real-space),
or use the link-variable Brillouin-zone construction in \Cref{subsec:bz-berry} (momentum-space).
\end{enumerate}

\textbf{Artifact controls (required for credibility).}
\begin{itemize}
\item \textbf{Gauge randomization test:} apply arbitrary local phase choices (rank-1) or local $U(n)$ frame rotations (rank-$n$) and
confirm that reported invariants (e.g.\ $\gamma_\Gamma$, $\mathrm{spec}(U_\Gamma)$) do not change.
\item \textbf{Resolution/convergence:} verify stability of invariants under refinement of $\Delta t$, lattice spacing, and loop sampling density.
\item \textbf{Adiabatic breakdown demonstration:} intentionally violate \Cref{eq:adiabatic-condition} and show that the extracted holonomy
becomes non-robust (strong dependence on sampling/path or visible subspace leakage), distinguishing geometry from beating.
\end{itemize}
\end{minipage}}
\end{center}

\subsection{Wilczek--Zee connection (non-Abelian)}

If the relevant polarization sector is an $n$-dimensional near-degenerate subspace, choose an orthonormal
frame $\mathbf{E}(x)\in\C^{N\times n}$ with columns $\lvert u_a(x)\rangle$ spanning the subspace, so that
$\mathbf{E}^\dagger\mathbf{E}=I_n$.
Define the Wilczek--Zee connection
\begin{equation}
  \mathcal{A}_\mu(x) := \ii\,\mathbf{E}(x)^\dagger\,\partial_\mu \mathbf{E}(x)\in\mathfrak{u}(n).
  \label{eq:wz-connection}
\end{equation}
A change of local frame $\mathbf{E}\mapsto \mathbf{E}\,U(x)$ with $U(x)\in U(n)$ yields the non-Abelian gauge law
\begin{equation}
  \mathcal{A}_\mu \mapsto U^\dagger \mathcal{A}_\mu U + \ii\,U^\dagger \partial_\mu U,
  \label{eq:wz-gauge}
\end{equation}
and the curvature (field strength) is
\begin{equation}
  \mathcal{F}_{\mu\nu}
  := \partial_\mu \mathcal{A}_\nu-\partial_\nu \mathcal{A}_\mu
     -\ii\,[\mathcal{A}_\mu,\mathcal{A}_\nu].
  \label{eq:wz-curvature}
\end{equation}

\subsection{\texorpdfstring{Axis-aligned triplet and prestress-driven $U(1)^3\rightarrow U(3)$ degeneracy}{Axis-aligned triplet and prestress-driven U(1)^3 -> U(3) degeneracy}}
\label{subsec:axis-triplet}

On the 6-neighbor axial-only cubic lattice (\Cref{subsec:lattice-substrate}) it is natural to organize the
per-wavepacket complex envelope as a three-component triplet $\Psi=(\psi_x,\psi_y,\psi_z)^\top\in\C^3$ aligned with the
lattice axes.
The role of the prestress parameter $\alpha$ (derivation convention: $\alpha = \mathrm{rest\_length}/\mathrm{spacing}$,
so $\alpha = 1$ is no prestress and $\alpha = 0$ is maximum prestress) is to control the \emph{eigenvalue spread} of the
3-channel Christoffel matrix $D(k)$ — not to introduce off-diagonal coupling between the axis components, which is
absent on this lattice at every $\alpha$.

\paragraph{Non-degenerate limit ($\alpha\to 1$, no prestress).}
The lattice energy is purely Hookean longitudinal, and $D(k)$ has maximally split eigenvalues with the transverse modes
at zero frequency. The three axis-aligned channels are dynamically decoupled: each component supports an independent
local phase choice and an independent rank-1 Berry connection $a_\mu^{(x)},a_\mu^{(y)},a_\mu^{(z)}$.
In this limit, it is meaningful (and gauge-invariant) to speak of ``the Berry phase of an axis.''
The carrier-triplet gauge group is closest to $U(1)^3$.

\paragraph{Degenerate limit ($\alpha\to 0$, maximum prestress).}
The $(1-\alpha)|\Delta u|^2$ geometric-stiffness term dominates and contributes equal stiffness to all components on
every bond. $D(k)\propto I$ becomes proportional to the identity at every $k$, and the three lateral channels are fully
degenerate. The transported subspace is genuinely 3-dimensional and the appropriate adiabatic-transport object is a
Wilczek--Zee connection $\mathcal{A}_\mu\in\mathfrak{u}(3)$ as in \Cref{eq:wz-connection}, decomposable into a trace
$U(1)$ part and a traceless $SU(3)$ part,
\begin{equation}
  \mathcal{A}_\mu
  = \frac{1}{3}\,\mathrm{tr}(\mathcal{A}_\mu)\,I_3
  + \sum_{a=1}^{8} A^a_\mu\,T^a,
\end{equation}
where $T^a$ are $SU(3)$ generators.
At the project's default operating point $\alpha=0.2$ the lattice sits close to this degenerate limit.

\paragraph{What the microscopic anisotropy actually provides.}
The $U(3)$ gauge group on a 3-dimensional complex envelope is a generic representation-theoretic fact about $\C^3$ and
does not require the lattice to be anisotropic.
What the cubic anisotropy does provide is the \emph{physically meaningful decomposition}: a preferred 3-dim subspace
(the lateral triplet, distinct from the embedding-amplitude direction $X^4$), a preferred axis-aligned basis within
it (which makes the eight $SU(3)$ traceless generators carry operational meaning, e.g.\ as ``colour'' labels on
solitons that occupy specific axis-aligned channels), and a natural identification of the $U(1)$ trace with the
electromagnetism-like sector.
Without the cubic anisotropy of the lab frame, no preferred basis would exist and the $SU(3)$ generators would be
arbitrary linear combinations with no operational handle.
This is the load-bearing role of microscopic anisotropy in the gauge-emergence story; see \texttt{paper/backbone.md} \#19
for the corresponding sector mapping.
That material anisotropy can carry genuine non-Abelian gauge structure in real space is not unique to the present
construction: \citet{Chen2019NonAbelianOptics} show that the off-diagonal permittivity/permeability of a wide class of
anisotropic optical media act as a pair of spin-dependent vector potentials, realizing a synthetic $SU(2)$ gauge field
in a homogeneous medium with an extractable Wilson loop and Zitterbewegung-like trajectories.
The present claim is the analogous statement for the mechanical polarization eigenbundle of the cubic brane, extended
to the triplet sector.

\subsection{Berry curvature on a discretized Brillouin zone (link-variable method)}
\label{subsec:bz-berry}

For a lattice substrate, Berry curvature can be computed directly on the discretized Brillouin zone using
link variables and plaquettes, closely analogous to lattice-QCD Berry-curvature constructions
\citep{Fukui2005ChernDiscretizedBZ,PuYamamoto2018BerryLatticeQCD,Yamamoto2016BerryPhaseLatticeQCD}.
Let $\lvert u(\mathbf{k})\rangle$ be a normalized eigenvector (or $\mathbf{E}(\mathbf{k})$ a frame of a degenerate
subspace) on a momentum lattice with spacing $\tilde a$.
In the Abelian case one defines links
\begin{equation}
  U_j(\mathbf{k})=\frac{\langle u(\mathbf{k})\,|\,u(\mathbf{k}+\tilde{\jmath})\rangle}
  {|\langle u(\mathbf{k})\,|\,u(\mathbf{k}+\tilde{\jmath})\rangle|},
\end{equation}
and a plaquette
\begin{equation}
  U_{ij}(\mathbf{k}) = U_i(\mathbf{k})U_j(\mathbf{k}+\tilde{\imath})
  U_i^\dagger(\mathbf{k}+\tilde{\jmath})U_j^\dagger(\mathbf{k}),
\end{equation}
whose argument approximates the integrated Berry curvature over the plaquette.
In the non-Abelian case one uses overlap matrices between frames and projects to $U(n)$ links, then forms the same
plaquette structure to obtain a discretized non-Abelian curvature.
This is a natural diagnostic bridge between the lattice-substrate hypothesis and QCD-motivated non-Abelian geometry.

\paragraph{Caveat: the construction is identity-by-construction on the real $D(k)$ eigenframe.}
The dynamical matrix $D(\mathbf{k})$ of the substrate is the Hessian of a real-valued spring energy, hence real and
symmetric for every $\mathbf{k}$. Real symmetric matrices admit real eigenvectors, and a real eigenvector bundle over a
simply-connected region of the Brillouin zone is topologically trivial: every plaquette holonomy reduces to the identity
and every continuous Berry curvature vanishes identically.
The link-variable construction above therefore yields zero curvature when applied to the eigenframe of the bare $D(k)$,
independent of $\alpha$ and the lattice details. (We have verified this numerically across a full $\alpha$ sweep and over
the entire $(k_x,k_y)$ plane; the holonomy is identity to machine precision at every plaquette and every $\alpha$.)
The meaningful Berry/WZ object in this project is not the bundle of real $D(k)$ eigenvectors but the
\emph{per-wavepacket complex envelope} $\Psi\in\C^3$ of \Cref{subsec:axis-triplet} together with its slow effective
Hamiltonian, whose eigenstates can carry genuine complex phase structure inherited from the rotating-wave reduction
$i\partial_t\Psi = H_{\mathrm{eff}}\Psi$.
Section~\ref{subsec:bz-berry} should therefore be read as a \emph{template diagnostic} that applies once the complex
envelope dynamics provides an intrinsically complex effective operator; applied to the lab-frame $D(k)$ alone it does
not separate $U(1)^3$ from $U(3)$.
\subsection{\texorpdfstring{Spinorial holonomy and the $4\pi$ aspect}{Spinorial holonomy and the 4pi aspect}}

For $n=2$ (a two-level polarization subspace), holonomy lives naturally in a double cover of $SO(3)$.
A closed loop in physical space that induces a $2\pi$ rotation of an internal polarization frame can therefore
produce a sign change of the transported state, while a $4\pi$ loop returns it to itself.
In the present theory this spinorial $4\pi$ aspect is attributed to Berry/Wilczek--Zee holonomy of an internal
polarization bundle, and is treated as part of the particle-sector structure developed in \Cref{sec:localized}.
